\section{Группы}

\begin{defn}{Множество с бинарной операцией}{}
    Множество \( M \), с заданным отображением \( M \times M \to M \) \( (a, b) \mapsto a \circ b \)

    Обозначается \( (M, \circ) \)
\end{defn}

\begin{defn}{Полугруппа}{}
    Множество с бинарной операцией \( (M, \circ) \) называется полугруппой,
    если эта операция ассоциативна:
    \[
        \forall a, b, c \in M \ (a \circ b) \circ c = a \circ (b \circ c)
    \]

    Обозначается \( (S, \circ) \)
\end{defn}

\begin{example}{Не всякая операция ассоциативна}{}
    \begin{itemize}
        \item
            \( M = \NN \) и \( a \circ b = a^b \) --- не ассоциативна:
            \[
                2^{(1^2)} \neq (2^1)^2
            \]
        \item
            \( M = \ZZ \) и \( a \circ b = a - b \) --- тоже не ассоциативна
    \end{itemize}
\end{example}

\begin{defn}{Моноид}{}
    Полугруппа \( (S, \circ) \) называется моноидом,
    если в ней существует нейтральный элемент:
    \[
        \exists e \in S: \ e \circ a = a \circ e = a
    \]
\end{defn}

\begin{example}{}{}
    Является ли \( (\NN, +) \) моноидом?
    Ответ зависит от географического положения,
    ведь в России не \( 0 \) не считают натуральным,
    а, например, во Франции, наоборот.
    Далее в курсе \( 0 \) не является натуральным.
\end{example}

\begin{remark}{}{}
    Если нейтральный элемент существует, то он единственный:
    \[
        e_1 = e_1 \circ e_2 = e_2
    \]
\end{remark}

\begin{defn}{Группа}{}
    Моноид \( (S, \circ) \) называется группой,
    если \( \forall a \in S \: \exists b \in S : a \circ b = b \circ a = e \)

    Обозначается \( (G, \circ) = (G, \cdot) = G \).

    Чаще всего используется мультипликативная запись: \\
    \( (a, b) \mapsto ab\), \( e = 1 \), обратный элемент обозначается \( a^{-1} \).
\end{defn}

\begin{exercise}{}{}
    Если обратный элемент существует, то он единственный.
\end{exercise}

\begin{defn}{Абелева группа}{}
    Группа \( G \) коммутативна (абелева), если \( \forall a, b \in G \ ab = ba \).

    В абелевых группах используется аддитивное обозначение: \\
    \( (a, b) \mapsto a + b\), \( e = 0 \), обратный элемент обозначается \( -a \).
\end{defn}

\begin{defn}{Порядок группы}{}
    Порядок группы \( G \) \( |G| \) --- число элементов в ней.

    \begin{itemize}
        \item
            \( |G| < \infty \Rightarrow G \) --- конечна
        \item
            \( |G| = \infty \Rightarrow G \) --- бесконечна
    \end{itemize}
\end{defn}

\begin{example}{Группы}{}
    \begin{enumerate}
        \item
            Числовые аддитивные:

            \( (\ZZ, +) \), \( (\QQ, +) \), \( (\RR, +) \), \( (\CC, +) \), \( (\ZZ_n, +) \)
        \item
            Числовые мультипликативные:

            \( (\QQ \setminus \{ 0 \}, \times) \),
            \( (\RR \setminus \{ 0 \}, \times) \),
            \( (\CC \setminus \{ 0 \}, \times) \),
            \( (\ZZ_p \setminus \{ 0 \}, \times) \)
        \item
            Группы матриц:

            \( GL_n (\RR) = \{ A \in M_n(\RR) \ \vline \ \det A \neq 0 \} \) --- общая линейная

            \( SL_n (\RR) = \{ A \in M_n(\RR) \ \vline \ \det A = 1 \} \) --- специальная линейная

            Есть еще куча разных групп: верхнетреугольные / нижнетреугольные / диагональные
            (конечно же, нужно взять только обратимые)
        \item
            Группы перестановок:

            \( S_n \), \( |S_n| = n! \)

            \( A_n \) --- четные перестановки, \( |A_n| = \frac{n!}{2} \)
    \end{enumerate}
\end{example}

\begin{exercise}{}{}
    \begin{itemize}
        \item
            \( S_n \) коммутативна \( \Leftrightarrow n \leq 2 \)
        \item
            \( A_n \) коммутативна \( \Leftrightarrow n \leq 3 \)
    \end{itemize}
\end{exercise}

\begin{defn}{Подгруппа}{}
    Подмножество \( H \) группы \( G \) называется подгруппой,
    если (определения эквивалентны):
    \begin{itemize}
        \item
            \begin{enumerate}
                \item
                    \( H \neq \varnothing \)
                \item
                    \( \forall a, b \in H \: ab^{-1} H \).
            \end{enumerate}
        \item
            \begin{enumerate}
                \item
                    \( e \in H\)
                \item
                    \( \forall a, b \in H \: ab \in H \)
                \item
                    \( \forall a \in H \: a^{-1} \in H \)
            \end{enumerate}
    \end{itemize}
\end{defn}

\begin{remark}{}{}
    В любой группе есть две подгруппы \( H = \{ e \} \) и \( H = G \)

    Они называются несобственными.
\end{remark}

\begin{prop}{}{}
    Любая подгруппа в \( (\ZZ, +) \) имеет вид \( k \ZZ \) (\( k \in \ZZ_{\geq 0} \))
\end{prop}

Очевидно, что \( k \ZZ \) --- подгруппа.

Пусть \( H \subseteq \ZZ \) --- подгруппа.

Если \( H = \{ 0 \} \), то \( k = 0 \).

Иначе в \( H \) есть ненулевой элемент.
Тогда в \( H \) есть положительный элемент.
Пусть \( k \) --- наименьший положительный элемент в \( H \).
Тогда \( k \ZZ \subseteq H \).
Пусть \( a \in H \).
Поделим с остатком: \( a = kq + r \), где \( 0 \leq r < k \).
\( r = a - kq \). \( a \in H, k \in H \Rightarrow r \in H \).
\( k \) наименьшее положительное число в \( H \), значит \( r = 0 \),
то есть \( a \in k \ZZ \).

\begin{defn}{Циклическая подгруппа}{}
    Пусть \( G \) --- группа, \( g \in G \).

    Циклической подгруппой,
    порожденной элементом \( g \) называется
    \[
        H = \langle g \rangle = \{ g^m, m \in \ZZ \}
    \]
\end{defn}

\begin{defn}{Порядок элемента}{}
    Пусть \( G \) --- группа, \( g \in G \)

    Порядком элемента \( g \) называется наименьшее натуральное число \( m \),
    для которого \( g^m = e \) или бесконечность, если такого \( m \) не существует.

    Обознается \( \ord (g) \)
\end{defn}

\begin{prop}{}{}
    \[
        | \langle g \rangle | = \ord (g)
    \]
\end{prop}
